\begin{abstract}
Este trabalho aborda a análise exploratória e a redução de dimensionalidade em dados de sensores de qualidade do ar, com foco em variáveis provenientes de sensores químicos de baixo custo. O objetivo é compreender o comportamento dos preditores e suas relações intrínsecas para subsidiar futuras etapas de modelagem e calibração de sensores. Utilizou-se o conjunto de dados \textit{Air Quality UCI} (2004--2005), composto por 9.358 registros horários de cinco sensores de óxidos metálicos (PT08.S1--S5) e três variáveis meteorológicas (temperatura, umidade relativa e umidade absoluta), além de medições de referência.

O método incluiu o pré-processamento das séries temporais (tratamento de ausentes, padronização e exclusão de variáveis com perda superior a 90\%), a análise estatística descritiva e a aplicação da Análise de Componentes Principais (PCA) implementada manualmente. Os resultados mostraram forte multicolinearidade entre sensores e correlações significativas com as variáveis de umidade, além de dois componentes principais capazes de explicar aproximadamente 77,5\% da variância total dos preditores.

Conclui-se que a PCA é uma ferramenta eficaz para reduzir a redundância dos dados, aumentar a interpretabilidade e evidenciar a influência combinada de fatores físico-químicos e ambientais nas respostas dos sensores, servindo como etapa essencial para análises preditivas e calibração de sensores de baixo custo.
\end{abstract}

\begin{IEEEkeywords}
Análise Exploratória de Dados, Qualidade do Ar, Sensores MOx, PCA, Séries Temporais, Multicolinearidade, Pré-processamento.
\end{IEEEkeywords}
